\begin{table}[htbp!]
	\centering
	\pgfplotstabletypeset[
	col sep=tab,
	header=true,
	zerofill,
	columns={0,1,2,3,4},
	columns/0/.style={
		column name={\text{\(U\) \([\si{\kilo\volt}]\)}},
		column type={S[table-format=1.1]},
		fixed zerofill,
		precision=1,
	},
	columns/1/.style={
		column name={\text{\(r_1\) \([\si{\milli\metre}]\)}},
		column type={S[table-format=1.2]},
		fixed zerofill,
		precision=2,
	},
	columns/2/.style={
		column name={\text{\(r_2\) \([\si{\milli\metre}]\)}},
		column type={S[table-format=1.2]},
		fixed zerofill,
		precision=2,
	},
	columns/3/.style={
		column name={\text{\(\lambda\) \([\si{\pico\metre}]\)}},
		column type={S[table-format=1.1]},
		fixed zerofill,
		precision=1,
	},
	columns/4/.style={
		column name={\text{\(\sigma_\lambda\) \([\si{\pico\metre}]\)}},
		column type={S[table-format=1.1]},
	},
	string type,
	every head row/.style={
		before row=\toprule,
		after row=\midrule,
	},
	every last row/.style={after row=\bottomrule}]{./Data/linear-data.tsv}
	\vspace{1.5pt}
	\caption{Datos del montaje de Tubo Plano.}
	\label{tab:linear-data}
\end{table}
